\documentclass[12pt,a4paper]{article}
%\usepackage[letterpaper, top=1.0in, bottom=1.0in, left=1.0in, right=1.0in]{geometry}
\usepackage[margin = 1.5cm]{geometry}

\usepackage{times}
\usepackage{bm}
%change section heading font size
\usepackage{titlesec}
\titleformat{\section}{\normalsize\bfseries}{\thesection}{1em}{}
\titleformat{\subsection}{\normalsize\bfseries}{\thesubsection}{1em}{}
\titleformat{\subsubsection}{\normalsize\bfseries}{\thesubsubsection}{1em}{}

%\usepackage{mathptmx}
\usepackage{cite,subfigure,epsfig,graphicx,epstopdf}
\usepackage[hyphens]{url}
\usepackage{amsmath,amssymb,amsfonts,amsthm,mathtools}
\usepackage{amsmath,wrapfig}
\usepackage{multirow}
\usepackage{multicol}
% \usepackage{arydshln}
\usepackage{array}
\usepackage{tabularx}
\usepackage{makecell}
\usepackage{lipsum}
% \usepackage[margin=1in]{geometry}
\usepackage{longtable}
\usepackage{setspace}
% \usepackage{algorithm}
% \usepackage{algpseudocode}
% \renewcommand{\algorithmicrequire}{\textbf{Input:}}
% \renewcommand{\algorithmicensure}{\textbf{Output:}}
\setlength{\fboxsep}{0pt}
\usepackage{enumitem}
\usepackage{hyperref}
% \usepackage{CJKutf8}
\usepackage{fontspec}
\usepackage{xeCJK}
% \AtBeginDvi{\input{zhwinfonts}}
\IfFontExistsTF{Times New Roman}
  {\setmainfont{Times New Roman}}
  {\setmainfont{TeX Gyre Termes}}
\IfFontExistsTF{Songti TC}
  {\setCJKmainfont{Songti TC}}
  {\IfFontExistsTF{Noto Serif CJK TC}
    {\setCJKmainfont{Noto Serif CJK TC}}
    {\setCJKmainfont{PingFang TC}}}
\hypersetup{colorlinks=true,linkcolor=black,urlcolor=blue,citecolor=black}

\usepackage{tikz}
\usepackage{pgfplots}
\usetikzlibrary{patterns}
\usepgfplotslibrary{groupplots}
\pgfplotsset{compat=1.12}

\newcommand{\RNum}[1]{\uppercase\expandafter{\romannumeral #1\relax}}
\newtheorem{definition}{\textbf{\emph{Definition}}}
\newcommand{\inlinesec}[1]{\noindent\textbf{#1.}}
\pagenumbering{arabic} \linespread{1}
\def\etal{\textit{et al}. }

\newcommand{\eg}{E_{\mathrm{g}}}
\newcommand{\ehull}{E_{\mathrm{hull}}}
\newcommand{\form}{x_{\mathrm{form}}}
\newcommand{\struc}{x_{\mathrm{struc}}}
\newcommand{\unc}{u}
\newcommand{\stab}{S_{\mathrm{stab}}}
\newcommand{\domainf}{d_{\mathrm{form}}}
\newcommand{\domains}{d_{\mathrm{struc}}}
\newcommand{\utility}{U}
\newcommand{\gapscore}{\rho_{\mathrm{gap}}}
\newcommand{\novel}{\rho_{\mathrm{novel}}}
\emergencystretch=2em
\setlength{\abovedisplayskip}{4pt}
\setlength{\belowdisplayskip}{4pt}
\setlength{\abovedisplayshortskip}{3pt}
\setlength{\belowdisplayshortskip}{3pt}

\begin{document}
\vspace{-60pt}
\begin{center}
    \textbf{CITY UNIVERSITY OF HONG KONG} \\
    \vspace{10pt}
    \underline{\textbf{Application Form for Zhongneng Tianyi Grant}} \\
\end{center}

\section{Name of the applicant, position and home Department, and the date of joining the University}

\begin{itemize}[leftmargin=15pt]
    \vspace{-8pt}
    \item Name of applicant: Prof. Chen MA (00473705)
    
    \vspace{-8pt}
    \item Position and affiliated department: Assistant Professor, Department of Computer Science, City University of Hong Kong
    
    \vspace{-8pt}
    \item Date of joining CityU: August 25, 2021
    
    \vspace{-8pt}
    \item Duration of the project: 2 years
\end{itemize}

\vspace{-20pt}
\section{Title of the Proposed Research}

\begin{itemize}
    \vspace{-8pt}
    \item English Project Title: AI-Powered Boron Nitride Material Exploration
    \vspace{-8pt}
    \item Chinese Project Title: 人工智能驅動的氮化硼材料探索
\end{itemize}

\vspace{-20pt}
\section{Abstract of the Project}
\vspace{-10pt}
This project aims to build an AI-powered boron nitride (BN) material exploration framework for prioritizing BN-related candidates in ultraviolet/wide-band-gap and dielectric/2D-support applications. The immediate challenge is not a lack of possible BN variants, but the absence of a \textbf{reliable decision layer} for sparse BN labels, formula-level extrapolation, and uncertain structure follow-up.

Our framework will make this decision process concrete: it will define a bounded BN design space, combine public 2D-material data with BN-centered diagnostics, screen formula-only candidates, and generate prototype-first structure hypotheses. Candidate decisions will use \textbf{target-window relevance, stability proxies, uncertainty, domain support, novelty, and action labels}.

Expected outputs include a curated dataset, benchmark report, ranked candidate table, validation-ready structure-handoff package, ranked candidate-review package for donor/technical discussion, and manuscript-ready analysis.

\vspace{-10pt}
\section{Key Objectives}

This project builds a BN-specific, uncertainty-aware, and application-oriented exploration pipeline. Our objectives are:

\begin{itemize}[leftmargin=*]
\item Construct a bounded BN-centered design space covering BN polymorphs, doped BN, BCN, BC$_2$N, SiBN, and chemically plausible B/N/C/Si-centered small-cell analogues.

\item Develop a two-stage AI screening process that first uses formula-only descriptors for broad candidate ranking and then applies structure-resolved checks after prototype-based structure hypotheses are available.

\item Create an uncertainty-aware decision layer for two application tracks: UV/wide-gap materials and dielectric support materials.

\item Deliver validation-ready evidence rather than unsupported discovery claims, including ranked families, action labels, structure files, uncertainty flags, and stakeholder-facing summaries.
\end{itemize}

\vspace{-20pt}
\section{Research Methodology}
\vspace{-10pt}

\subsection{Background of Research}
\label{sec:background}
BN is a focused yet scientifically rich materials family. h-BN and c-BN anchor the known design space, while B-C-N, BC$_2$N, Si-B-N, doped BN, and controlled substitutions define a broader but still bounded exploration region. BN-based materials are relevant to ultraviolet emission, wide-band-gap electronics, insulating layers, and 2D heterostructure supports~\cite{watanabe2004_hbn_uv,cassabois2016_hbn_indirect,dean2010_hbn_graphene}. The core research question is therefore not whether AI can predict every material, but \textbf{which BN-related formulas and structures should be investigated next, and with what uncertainty}.

Preliminary implementation has established PoC-level modules for public 2D-material data ingestion, grouped-by-formula evaluation, BN formula/family diagnostics, candidate-compatible screening models, ranking uncertainty layers, and first-pass structure follow-up artifacts. These results support \textbf{BN-themed candidate prioritization and follow-up planning}, not direct claims of BN material discovery or synthesis.

\vspace{-10pt}
\subsection{Research on AI-Powered BN Material Exploration}
\label{sec:bn-ai-exploration}
The proposed framework treats materials exploration as a decision pipeline rather than a single prediction task. The bounded candidate space will be represented as
\begin{equation}
  \mathcal{C}_{\mathrm{BN}}=
  \mathcal{C}_{\mathrm{core}}\cup\mathcal{C}_{\mathrm{ext}}\cup\mathcal{C}_{\mathrm{explore}},
  \quad
  \mathcal{C}_{\mathrm{explore}}=
  \{c:q(c)=0,\;N_{\mathrm{atom}}(c)\le 12,\;\pi(c)\in\Pi_{\mathrm{BN}}\}.
\end{equation}
Here, $\mathcal{C}_{\mathrm{core}}$ covers BN polymorphs and doped BN, $\mathcal{C}_{\mathrm{ext}}$ covers BCN, BC$_2$N and SiBN motifs, and $\mathcal{C}_{\mathrm{explore}}$ is restricted to charge-neutral small-cell substitutions that can be mapped to known BN-family structural motifs. This boundary follows the project goal of being general enough for exploration but not so open-ended that the search becomes chemically meaningless.

\vspace{-10pt}
\subsubsection{Subtask on BN-Centered Data Construction and Diagnostics}
\label{subsubsec:data-diagnostics}
\vspace{-5pt}
The first subtask is to construct a provenance-aware BN data layer from public computational resources such as 2DMatPedia, JARVIS, and the Materials Project~\cite{zhou2019_2dmatpedia,choudhary2020_jarvis,jain2013_mp}. Benchmarking, featurization, and 2D-material cross-checks will use Matbench, matminer, and C2DB where appropriate~\cite{dunn2020_matbench,ward2018_matminer,haastrup2018_c2db}. Each record is represented as
\begin{equation}
  r_i=(c_i,s_i,\mathbf{y}_i,p_i),
  \qquad
  p_i=(\mathrm{source},\mathrm{method},\mathrm{version},\mathrm{provenance}),
\end{equation}
where $c_i$ is composition, $s_i$ is optional structure, and $\mathbf{y}_i$ stores available target properties. Instead of merging all labels as identical ground truth, incompatible labels will be retained with source indicators. This is essential for a BN-centered project because current BN labels are sparse and unevenly distributed.

The diagnostics will explicitly separate general 2D-material performance from BN-centered generalization. The existing codebase already distinguishes grouped-by-formula evaluation from BN formula holdout, BN family holdout, and BN/non-BN stratified error. Success criteria will include grouped-by-formula MAE/RMSE for band-gap prediction, BN-family holdout error, uncertainty coverage on held-out candidates, and the number of candidates promoted to structure-level review. A candidate family will be promoted only if ensemble ranking remains stable under objective-weight perturbation and the structure-level recheck does not reverse the target-window or stability decision.

\vspace{-10pt}
\subsubsection{Subtask on Formula-Only Screening and Uncertainty-Aware Ranking}
\label{subsubsec:screening-ranking}
\vspace{-5pt}
The second subtask is to develop a two-stage screening route. Formula-only candidates will first be evaluated using composition descriptors and composition-based models:
\begin{equation}
  f_{\theta}^{(1)}:\form\rightarrow(\hat{\eg},\hat{\stab},\hat{\unc}),
  \quad
  f_{\theta}^{(2)}:(\form,\struc)\rightarrow
  (\hat{\eg},\hat{\ehull},\hat{\boldsymbol{\epsilon}},\hat{\unc}).
\end{equation}
The first function supports wide candidate screening; the second is used only after structure hypotheses are available. This distinction matches the current project code, where the best overall model and the best candidate-compatible screening model are not assumed to be the same. At the formula-only stage, stability and application relevance will be treated as composition-level proxies only. Structure-dependent quantities such as energy above hull and dielectric tensors will be reported only after prototype-based structure hypotheses and structure-resolved checks are available.

The band-gap term will be treated as a target-window score rather than a monotonic objective:
\begin{equation}
  \gapscore(c)=
  \exp\!\left[-\frac{(\hat{\eg}(c)-E^*_{\mathrm{g,app}})^2}{2\sigma_{\mathrm{app}}^2}\right]\eta_{\mathrm{direct}}(c).
\end{equation}
At formula-only stage, directness is unknown unless explicitly labelled; it is applied conservatively after structure-resolved scoring. Application scores for UV/wide-band-gap and dielectric/2D-support tracks are proxy scores derived from computed properties, dimensionality metadata, and rule-based target criteria. Predictive uncertainty will be estimated by model ensembles and validation-error calibration~\cite{gal2016_dropout,lakshminarayanan2017_ensembles}:
\begin{equation}
  \hat{\unc}_{\mathrm{cal}}(c)=
  \alpha\,\mathrm{Std}_{m\in\mathcal{M}}[f_m(c)]
  +\beta\,\widehat{\mathrm{err}}_{\mathrm{val}}(c).
\end{equation}

Candidate utility combines performance, reliability, domain support, and novelty:
\begin{equation}
  \begin{aligned}
  \utility_t(c,s)=&\;w_1\rho_t(c,s)+w_2\hat{\stab}(c,s)-w_3\hat{\unc}(c,s)\\
  &-w_4 d_t(c,s)+w_5\novel(c),
  \end{aligned}
  \quad t\in\{\mathrm{UV/WBG},\mathrm{dielectric/2D}\}.
\end{equation}
Domain distance is formula-level before structure handoff and structure-level afterwards: $d_t=\domainf(c)$ or $\domains(c,s)$. The output action label will be one of \textit{control}, \textit{priority}, \textit{explore}, or \textit{hold}. Provisional rules are auditable rather than fixed as final constants: priority means top-quartile utility, calibrated uncertainty below the 75th percentile of validation error, non-outlier domain distance, and no failed stability or charge sanity flag; explore means high novelty or top-quartile utility with above-threshold uncertainty; control means known BN-family references; hold means target-window failure, outlier domain distance, or excessive uncertainty. The UV/WBG track will initially target a 4.5--6.5 eV gap window, with direct-gap evidence treated as a bonus after structure handoff. The dielectric/2D-support track will prioritize insulating candidates with large predicted gap, 2D provenance or plausible layered prototype, low uncertainty, and no stability sanity failure. Ranking reliability will be measured by calibration error, ensemble rank stability, and top-$K$ hit rate after structure-level rechecking~\cite{lookman2019_active}.

\vspace{-10pt}
\subsubsection{Subtask on Structure Hypothesis Handoff and Validation}
\label{subsubsec:structure-validation}
\vspace{-5pt}
The third subtask connects formula-level ranking to structure-level evidence through a staged protocol. Selected formulas will first be converted into candidate structures through prototype reuse and database-derived structure transfer; controlled substitution will be used when a close BN-family prototype exists but the exact formula is absent. Generative approaches such as CDVAE will be considered only when no plausible prototype-based route is available~\cite{xie2022_cdvae}. Structure-resolved graph or potential models such as CGCNN, MEGNet, ALIGNN, and CHGNet will be used only after suitable structures pass sanity checks~\cite{xie2018_cgcnn,chen2019_megnet,choudhary2021_alignn,deng2023_chgnet}.

Materials-side validation will include charge-neutrality checks, structural sanity checks, relaxation-status checks, energy-above-hull or formation-energy proxies, property recheck, and comparison with known BN-family references. For bulk-like entries, the stability proxy will focus on formation energy or energy-above-hull:
\begin{equation}
  \ehull(c,s)=e(c,s)-
  \min_{\lambda_q\ge0}\sum_q\lambda_q e(q),
  \quad \sum_q\lambda_q=1,\;\sum_q\lambda_q\mathbf{n}_q=\mathbf{n}_c.
\end{equation}
For 2D candidates, dimensionality, database provenance, and available 2D stability indicators will also be considered. The purpose is to prepare a validation-ready handoff package, not to claim experimental synthesis.

\vspace{-10pt}
\subsection{Significance, Expected Outcomes, and Major Deliverables}
\label{sec:significance}
The significance of this project is a BN-specific AI decision pipeline: a curated data layer, family-aware diagnostics, calibrated uncertainty, two-track candidate ranking, and prototype-first structure handoff. It supports the PI in deciding which BN-related candidate families deserve additional computational or stakeholder review.

Expected deliverables follow the output chain
\begin{equation}
  \mathrm{BN\ dataset}\rightarrow\mathrm{benchmarked\ models}\rightarrow
  \mathrm{ranked\ candidates}\rightarrow\mathrm{structure\ handoff}\rightarrow
  \mathrm{technical\ report}.
\end{equation}
The package will include a BN-centered dataset with formula, structure, target-property, dimensionality, and provenance fields; a benchmark report under random, grouped, and BN-family-aware splits; a ranked table of at least 30 candidate families with predicted properties, uncertainty, domain support, novelty score, application label, and recommended action; a structure hypothesis package for 5--10 priority candidates with CIF/POSCAR files and validation notes; and manuscript-ready analysis with a lightweight demo interface.

The initial phase is computational and data-driven. No human subjects, living animals, wet-lab synthesis, biological samples, or radiation work will be conducted. If later experimental synthesis or regulated laboratory procedures are introduced, the responsible experimental party will obtain the required approvals before that work begins.

\vspace{-10pt}
\section{Requested Funding and Budget}
\vspace{-10pt}
The total funding requested is \$330,000 HKD, which is broken down as follows:
\vspace{-6pt}
\begin{center}
\footnotesize
\begin{tabular}{|>{\raggedright\arraybackslash}p{0.15\linewidth}|>{\raggedright\arraybackslash}p{0.56\linewidth}|>{\raggedright\arraybackslash}p{0.13\linewidth}|}
\hline
\textbf{Budget item} & \textbf{Detail breakdown and justification} & \textbf{Estimated cost} \\
\hline
Staffing & Recruiting one MPhil student & 310,000 HKD\\
& Justification: 19,100 HKD $\times$ 12 months $\times$ 1 person = 229,200 HKD & \\
& Recruiting one RA & \\
& Justification: 13,466.6 HKD $\times$ 6 months $\times$ 1 person = 80,800 HKD & \\
\hline
Conference & The research outcomes of this project are expected to be published at & 20,000 HKD \\
& top-tier conferences, e.g., ACM SIGKDD. & \\
\hline
\multicolumn{2}{|r|}{\textbf{Total}}   & \textbf{330,000 HKD} \\
\hline
\end{tabular}
\normalsize
\vspace{-12pt}
\end{center}

\let\oldbibliography\thebibliography
\renewcommand{\thebibliography}[1]{%
  \oldbibliography{#1}%
  \setlength{\itemsep}{-5pt}%
}
%%
\bibliographystyle{IEEEtran}
\bibliography{ai_for_bn_research_plan_v18}

\end{document}
